\documentclass{article}

% The root Makefile supplies the path to the unmodified ICML style files.
% Shared, topic-agnostic preamble for all three papers.
% Keep formatting changes here conservative: ICML forbids modifications that
% create a space advantage.
\usepackage{microtype}
\usepackage{graphicx}
\usepackage{subcaption}
\usepackage{booktabs}
\usepackage{hyperref}

% Improve interoperability between hyperref and algorithmic.
\newcommand{\theHalgorithm}{\arabic{algorithm}}

% Review mode is intentionally the default. Change this to
% \usepackage[accepted]{icml2026} only after acceptance.
\usepackage{icml2026}

\usepackage{amsmath}
\usepackage{amssymb}
\usepackage{mathtools}
\usepackage{amsthm}
\usepackage[capitalize,noabbrev]{cleveref}

\theoremstyle{plain}
\newtheorem{theorem}{Theorem}[section]
\newtheorem{proposition}[theorem]{Proposition}
\newtheorem{lemma}[theorem]{Lemma}
\newtheorem{corollary}[theorem]{Corollary}
\theoremstyle{definition}
\newtheorem{definition}[theorem]{Definition}
\newtheorem{assumption}[theorem]{Assumption}
\theoremstyle{remark}
\newtheorem{remark}[theorem]{Remark}

\graphicspath{{figures/}}

% Visible by design: no template should be mistaken for a completed paper.
\newcommand{\placeholder}[1]{%
  \par\noindent\textit{Template note: #1}\par%
}


\icmltitlerunning{Adaptive-Mixture Policy Gradients for Length-Constrained RLVR}

\begin{document}

\twocolumn[
  \icmltitle{Adaptive-Mixture Augmented Policy Gradients for
  Expected-Length-Constrained Reinforcement Learning with Verifiable Rewards}

  % Author details remain hidden in review mode. Replace these placeholders
  % only when preparing the camera-ready version.
  \begin{icmlauthorlist}
    \icmlauthor{Anonymous Author(s)}{anon}
  \end{icmlauthorlist}
  \icmlaffiliation{anon}{Anonymous Institution}
  \icmlcorrespondingauthor{Anonymous Author}{anonymous@example.com}

  \icmlkeywords{reinforcement learning, verifiable rewards, importance
  sampling, post-training, language models}

  \vskip 0.3in
]

\printAffiliationsAndNotice{}

\begin{abstract}
Reinforcement learning with verifiable rewards spends most of its compute on
fresh language-model trajectories, while ordinary replay is statistically
problematic because sequence likelihood ratios become unstable and previously
consumed samples depend on the current policy. We study a narrower alternative:
at each frozen policy, combine fresh on-policy trajectories with newly
generated, single-use trajectories from a known stale actor and an adaptive
prompt distribution. A current-policy mixture component bounds the exact joint
importance ratio, and frozen reward and length models supply an augmented
control variate. Under explicit sampling and decoding assumptions, the
resulting one-update gradient estimator is conditionally
unbiased. A Lagrange multiplier constrains expected response length instead of
silently normalizing the objective per token. This
methods-first draft gives the estimator audit and matched-budget training
protocol, but reports no performance gain before those experiments are run.
\end{abstract}

\section{Introduction}
\label{sec:introduction}

Language-model reinforcement learning with verifiable rewards (RLVR) can train
reasoning behavior from executable tests or exact answer checks. Its simplicity
does not remove the cost of generation: a policy update may require many long
trajectories, most of which are discarded after one use. Critic-free
descendants of REINFORCE \cite{williams1992reinforce}, such as RLOO
\cite{ahmadian2024back} and GRPO
\cite{shao2024deepseekmath} reduce value-model overhead, while large-scale
systems add clipping, prompt filtering, and token-level aggregation
\cite{yu2025dapo}. These choices improve engineering behavior but change which
prompt, rollout, or token population receives weight.

Reusing trajectories appears attractive, yet ``off-policy'' covers distinct
statistical settings. A replay-buffer trajectory may have selected itself into
the buffer through its reward and may already have changed the current policy.
Recomputing its likelihood does not restore independence. Sequence-level
importance ratios can have extreme variance; token-level ratios do not
represent the sampling probability of a sequence-rewarded trajectory.
Practical off-policy methods therefore use tapering, buffers, or adaptive
clipping \cite{leroux2025topr,wan2026buffer,xi2026bapo}, knowingly trading bias
for stability.

We ask whether a controlled source of stale data can improve gradient
efficiency without making an inaccurate replay-unbiasedness claim. At iteration
$t$, a live stale actor produces \emph{new}, never-consumed
trajectories from a logged policy and prompt distribution. These are combined
with a positive fraction of current-policy trajectories. An augmented
importance estimator uses a frozen outcome predictor to reduce residual
variance, then takes exactly one policy update and consumes the data.
This does not reuse rollout FLOPs: its plausible benefit is lower
residual-gradient variance per total generation and scoring cost when the
alternative proposal reaches informative regions. Current and stale actors may
generate in parallel while the learner remains frozen. Learner--actor
asynchrony is a separate non-theorem practical variant, not a claimed benefit
of the core algorithm.

Our contributions are deliberately scoped:
\begin{itemize}
  \item We define the target joint trajectory law, including the prompt and
  decoding distributions, and derive a stratified adaptive-mixture augmented
  policy-gradient estimator with bounded exact density ratio.
  \item We state sufficient conditions for conditional unbiasedness and the
  precise two-route bias identity. The result excludes ordinary replay,
  multiple epochs, clipping, and group standardization.
  \item We couple the estimator to an explicit expected-token constraint and
  provide logging and diagnostics that can falsify its assumptions.
  \item We specify a frozen-policy estimator audit and a one-update-per-batch
  matched-budget study across model families and verifier domains.
\end{itemize}

\paragraph{Evidence status.}
No downstream or efficiency result is asserted in this draft. The theorem
concerns one raw gradient estimate at a frozen policy; it neither guarantees a
better optimizer step nor policy improvement.

% Camera-ready only: if the work evaluates technology developed by an author's
% employer or has another substantive financial conflict, place the required
% "Conflict of Interest Disclosure" paragraph at the end of this section.

\section{Related Work}
\label{sec:related}

\paragraph{On-policy and group-relative RL.}
Careful PPO reproduction shows that reward whitening, KL treatment, clipping,
and response length are implementation-sensitive
\cite{huang2024nplus}. RLOO gives a simple leave-one-out REINFORCE baseline
\cite{ahmadian2024back}; GRPO standardizes rewards within prompt groups
\cite{shao2024deepseekmath}; Dr.\ GRPO questions group-standard-deviation and
response-length normalization \cite{liu2025critical}; and DAPO adds dynamic
sampling, asymmetric clipping, and token-level aggregation
\cite{yu2025dapo}. GSPO uses a length-normalized geometric mean of token ratios
with response-level clipping \cite{zheng2025gspo}. Recent analyses expose
negative-gradient pathologies
\cite{deng2025negative}, U-statistic structure \cite{zhou2026ustatistic}, and
the interaction between clipping, entropy, and misaligned rewards
\cite{chen2026exploration}. Our estimator intentionally omits group
normalization and PPO clipping so its target expectation remains explicit.

\paragraph{Allocation, baselines, and length.}
Adaptive rollout allocation predicts prompt difficulty and assigns rollout
counts under a fixed budget \cite{nguyen2026adaptive}. Cross-prompt shrinkage
baselines reduce variance while preserving a policy-gradient target
\cite{zeng2026shrinkage}. Group-relative reward rescaling addresses length
inflation \cite{li2026length}. We instead represent length as a constrained
cost and adapt the stale-source prompt proposal; these components are not
claimed to subsume those methods.
Discounted Beta--Bernoulli reward estimation reuses historical reward
statistics to reduce group-estimator error while explicitly accepting bias
\cite{kim2026dbb}. It is therefore a required bias--variance comparator:
reusing summary statistics is distinct from our density-corrected use of newly
drawn stale-policy trajectories.

\paragraph{Off-policy stability and reward validity.}
TOPR treats positive examples with an SFT/log-ratio-like term and truncates
importance weighting for negative examples \cite{leroux2025topr};
buffer-based and balanced policy optimization improve
practical reuse \cite{wan2026buffer,xi2026bapo}. Such stabilized estimators may
be preferable even when biased. Separately, spurious-reward controls show that
benchmark gains need not demonstrate learning from the intended verifier
\cite{shao2026spurious}. The protocol includes those controls and treats exact statistical
assumptions, not empirical stability alone, as the object of the first study.
Direct replay studies optimize the generation--staleness--diversity trade-off
\cite{arnal2026replay}, while POPO prioritizes effective recent groups and uses
decoupled off-policy correction \cite{mao2026popo}. Both are stronger practical
comparators than naive reuse. Their empirical utility does not imply the
conditional unbiasedness claimed for our newly sampled, single-use design.
Asynchronous RLHF separates generation from learning on lagged policies
\cite{noukhovitch2025async}; it motivates the stale-actor systems layout but
does not provide our conditional estimator identity. Statistically, the design
is a restricted stratified multiple-importance sampler with augmentation.
Classical work optimizes mixture allocation and control variates jointly
\cite{he2014optimal}; our pilot deliberately studies a small fixed grid rather
than claiming a new general MIS principle or optimal allocation.

PrefixRL and PROS instead reuse an old prefix while generating a new
continuation \cite{setlur2026prefixrl,huang2026pros}; they change the prompt or
conditional objective and avoid full-trajectory ratios. They are practical
compute comparators, not instances of \cref{eq:estimator}. At the systems
layer, HybridFlow decouples RL control and distributed computation
\cite{sheng2025hybridflow}; our implementation follows that separation so
rollout, cross-policy scoring, and update costs remain individually measurable.
In its exact-acceptance setting, SPEC-RL uses old trajectories as speculative
drafts while preserving the current-policy rollout law; its lenient setting
does not have that exact claim \cite{liu2025specrl}. MinPRO instead replaces exact causal-prefix ratios by a
stabilized minimum-ratio surrogate under policy lag \cite{lei2026minpro}.
These methods target complementary systems and biased-stability trade-offs.

\section{Formal Setting}
\label{sec:setting}

Let $Z=(X,Y)$ contain prompt $X$ and a complete generated sequence $Y$,
including EOS or a declared truncation event. The target prompt distribution is
$D$, a versioned finite prompt multiset with known masses $d(x)$, and the exact rollout distribution is
$\pi_\theta(y\mid x)$. Thus
\begin{equation}
P_\theta(z)=d(x)\pi_\theta(y\mid x).
\label{eq:target-law}
\end{equation}
The rollout decoder is part of $\pi_\theta$. Temperature, truncation, and any
rejection rule must be represented by its normalized sampling probability;
raw base-model logits are not automatically the rollout density. The
theorem-facing implementation uses temperature-softmax sampling without
top-$k$, top-$p$, repetition penalties, minimum-length masks, or forced tokens,
with fixed maximum length and EOS included in the outcome.
The allowed vocabulary $\mathcal V_\ell$ is a deterministic,
parameter-independent function of the tokenizer and prefix. For temperature
$T_d$ the logged
per-token mass is
\begin{equation}
p_\theta(y_\ell\mid y_{<\ell},x)=
\frac{\mathbf 1\{y_\ell\in\mathcal V_\ell\}
\exp(z_{\theta,\ell,y_\ell}/T_d)}
{\sum_{v\in\mathcal V_\ell}\exp(z_{\theta,\ell,v}/T_d)}.
\label{eq:decoder-pmf}
\end{equation}
The sequence mass is the product through EOS; no-EOS sequences at the fixed
limit carry their product mass and a deterministic truncation status. An
empirical conformance test for both current and stale sampler--scorer pairs
must pass before estimator audits.

The verifier deterministically maps every complete $Z$, including parse,
timeout, sandbox-error, and truncation statuses, to bounded reward $R(Z)$; no
outcome is retried or dropped. Interactive or stochastic environments are
outside the core formulation. Generation incurs $C(Y)=|Y|$ response tokens.
For fixed expected-response budget $\kappa$, we optimize the Lagrangian
\begin{equation}
\mathcal L(\theta,\lambda)=
\mathbb E_{P_\theta}[R]-
\lambda\{\mathbb E_{P_\theta}[C]-\kappa\},
\quad \lambda\ge0.
\label{eq:lagrangian}
\end{equation}
The constrained problem is
$\max_\theta\min_{\lambda\ge0}\mathcal L(\theta,\lambda)$ with
$\mathbb E_{P_\theta}[C]\le\kappa$. The stochastic dual update below is a
constraint controller, not a finite-sample feasibility guarantee.
Define $U_\lambda=R-\lambda C$ and the raw sequence score
\begin{align}
s_\theta(Z)
&=\left.\nabla_\vartheta
\log\pi_\vartheta^{\mathrm{decoder}}(Y\mid X)
\right|_{\vartheta=\theta}\nonumber\\
&=\sum_{\ell=1}^{|Y|}
\left.\nabla_\vartheta\log
\pi_\vartheta^{\mathrm{decoder}}(Y_\ell\mid X,Y_{<\ell})
\right|_{\vartheta=\theta}.
\label{eq:score}
\end{align}
We do not divide \cref{eq:score} by realized length. Under standard
score-function conditions,
$\nabla_\theta\mathcal L=\mathbb E_{P_\theta}[U_\lambda s_\theta]$.

\section{Method}
\label{sec:method}

Condition on the pre-update history $\mathcal F_t$. The current policy
$\theta_t$, dual variable, nuisance models, and all sampling distributions are
frozen before trajectories are drawn. A logged adaptive prompt distribution
$Q_t$ and a stale behavior policy $\mu_t$ define
\begin{equation}
G_t(z)=q_t(x)\mu_t(y\mid x).
\end{equation}
$G_t$ is generated only after all iteration-$t$ choices are frozen. A queue
whose realized contents can affect $\theta_t$, $q_t$, $\alpha$, or nuisance
models is not conditionally fresh and does not satisfy the theorem, even if
its trajectories have never been optimized on.

For fixed $\alpha\in(0,1]$, define the joint mixture
\begin{equation}
M_t(z)=\alpha P_{\theta_t}(z)+(1-\alpha)G_t(z)
\end{equation}
and exact detached weight
\begin{equation}
w_t(z)=\frac{P_{\theta_t}(z)}{M_t(z)}
\le \frac{1}{\alpha}.
\label{eq:weight}
\end{equation}
The prompt density is part of \cref{eq:weight}. A conditional policy ratio is
valid only when $Q_t=D$.
Defensive mixtures and augmented importance estimators are established
statistical tools \cite{owen2000safe,dudik2011doubly}; our contribution is
their restricted, auditable instantiation for fresh stale-actor RLVR, not the
general identities themselves.

\subsection{Learning Objective}

Let $h_t^R(z)$ and $h_t^C(z)$ be frozen action-conditional predictors of reward
and length, and set $h_t=h_t^R-\lambda_th_t^C$. This separation allows the
dual variable to change without retraining the two predictors. Let $b_t(x)$ be
a frozen prompt-only baseline. In stale-source conditions, fix integers
$n_P,n_G\ge1$ and mixture mass $\alpha\in(0,1)$ independently before sampling;
write $\rho=n_P/(n_P+n_G)$ for the allocation fraction. Draw
$Z_i^P\sim P_{\theta_t}$ and $Z_j^G\sim G_t$. For $\alpha<1$, deterministic
stratification gives
\begin{align}
\widehat g_t={}&
\frac1{n_P}\sum_{i=1}^{n_P}
\Big[(h_i-b_i)s_i+
\alpha w_i(U_i-h_i)s_i\Big]\nonumber\\
&+\frac{1-\alpha}{n_G}\sum_{j=1}^{n_G}
w_j(U_j-h_j)s_j.
\label{eq:estimator}
\end{align}
The first line uses current-policy samples and the second uses stale-source
samples. All $w,h,b$ are stop-gradient quantities. The $\alpha=1$ limit omits
the stale-source sum and sets $n_G=0$; it reduces exactly to ordinary on-policy
$(U-b)s$. The outcome predictor is critic-like infrastructure; we do not call
the method critic-free. The exploratory pilot restricts $\rho=\alpha$ to keep
64 total trajectories; the identity itself does not.

For conditionally IID trajectories within each of two independent strata define
$A=(h-b)s+\alpha w(U-h)s$ under $P$ and
$B=w(U-h)s$ under $G$. Conditional on the frozen history,
\begin{equation}
\operatorname{Cov}(\widehat g_t)=
\frac{\operatorname{Cov}_P(A)}{n_P}+
\frac{(1-\alpha)^2\operatorname{Cov}_G(B)}{n_G}.
\label{eq:variance}
\end{equation}
This expression includes covariance between direct and residual terms inside
$A$. It gives no universal advantage over on-policy sampling; the
frozen-policy audit chooses $\alpha$ from a preregistered grid using projected
MSE times measured unit cost, then freezes it for confirmation. Nuisance fitting,
$\alpha$ selection, and confirmation use three disjoint prompt partitions; the
projection seed, dimension, reference sample, and confirmation replications
are frozen before selection.
Before these partitions are formed, every prompt is tied to a pinned source
revision, export-byte hash, stable upstream identifier, raw-row hash, and
versioned normalized-content hash. The downstream partitioner assigns every
retained duplicate cluster wholly to fit, selection, or confirmation; any
later admitted hidden evaluator is a separate registered role. Semantic
overlap with such evaluators is a blocking audit rather than being inferred
from normalized-hash inequality. A hash-bound preflight rejects named
PII/secret patterns and retained character-5-gram near-duplicates, but this
lexical heuristic does not replace reviewed multilingual, paraphrase, or
hidden-evaluator leakage analysis.

$h_t^R,h_t^C,b_t$, and $q_t$ are trained only from data available before iteration
$t$, then frozen. The sums below range over the complete finite prompt
multiset, and prompts are sampled with replacement. A convenient adaptive proposal is
\begin{align}
a_t(x)&=d(x)\sqrt{\widehat v_t(x)/\widehat c_t(x)},\\
Z_t^q&=\sum_{x'}a_t(x'),\\
\overline q_t(x)&=a_t(x)/Z_t^q,\\
q_t(x)&=\epsilon d(x)+(1-\epsilon)\overline q_t(x).
\end{align}
where $\widehat c_t(x)>0$ predicts generation cost and $\widehat v_t(x)$ uses
prior-iteration random projections to predict stale-source residual-gradient
second moment. Because $q$ also changes the mixture weight and its covariance
with the direct term, this square-root rule is a heuristic, not an optimum for
\cref{eq:estimator}. Incorrect predictors may reduce efficiency but do not
alter the density-route expectation when the logged distribution and ratio are
exact.
We impose $\widehat v_t(x)\ge v_{\min}>0$ before normalization; equivalently,
an implementation may set $\overline q_t=D$ whenever $Z_t^q=0$.

\subsection{Feedback and Credit Assignment}

The main study uses deterministic programmatic verifiers: normalized answer
checking for mathematics and isolated executable tests for code. Parser
failures, timeouts, sandbox errors, and truncations are separate logged outcomes
with predeclared scalar rewards rather than silently dropped samples. Hidden
tests are disjoint from the authors' SFT, nuisance-model, and prompt-allocation
data; unknown base-model pretraining contamination remains a limitation.
Source and prompt draws are assigned before distributed generation. A valid
sequence that reaches any declared verifier status is retained regardless of
completion time. An infrastructure failure before a scored sequence exists
aborts the complete batch, charges incurred cost, and performs no update;
workers are never refilled selectively.
This prevents selective refill but does not make completion ignorable. The
theorem-facing audit therefore requires every preassigned draw to reach a
scored outcome almost surely. If any infrastructure failure occurs, that
replication is retained as an operational failure, never retried, and excluded
from the estimator claim; H1 cannot pass from a success-conditioned subset.

Terminal $U_\lambda$ multiplies the unnormalized sequence score in
\cref{eq:score}; the theorem does not assign token-level causal credit. Reward
shaping, group centering, group standardization, token-average losses, ratio
clipping, and gradient clipping all change either the estimand or estimator and
are excluded from the theorem-facing condition. They may be evaluated as
clearly labeled practical variants. Shuffled, random, weakly related, and
anti-correlated verifier controls test whether improvement follows the intended
signal rather than pretrained behavior or evaluation leakage.

\subsection{Algorithm and Complexity}

\begin{algorithm}[t]
\caption{One Adaptive-Mixture Update}
\label{alg:am}
\begin{algorithmic}[1]
\STATE Freeze $\theta_t,\lambda_t,h_t^R,h_t^C,b_t,q_t,\mu_t,$ and $\alpha$.
\STATE Draw $n_P$ complete trajectories from $D\pi_{\theta_t}$.
\STATE Draw $n_G$ new, single-use trajectories from $Q_t\mu_t$.
\STATE Log exact prompt, source, decoder, and sequence probabilities.
\STATE Evaluate verifier reward and token cost; form $U=R-\lambda_tC$.
\STATE Compute detached $w=P_{\theta_t}/M_t$ and \cref{eq:estimator}.
\STATE Finish accumulation at $\theta_t$; call \texttt{optimizer.step()} once.
\STATE Consume every trajectory; no microbatch update or second epoch.
\STATE Update $\lambda$ from fresh current-policy token costs.
\STATE Add consumed data only to future nuisance-model training.
\end{algorithmic}
\end{algorithm}

Relative to on-policy REINFORCE, the method adds stale-actor generation, current
and stale sequence log-probabilities under both policies, and nuisance-model
inference. The cost ledger includes all generated tokens, verifier calls,
likelihood evaluations, nuisance training, and policy forwards/backwards.
Storage is linear in the single-use queue. A practical repeated-replay variant
can reduce rollout cost, but is outside the theorem and is reported as biased.

\section{Analysis}
\label{sec:analysis}

\begin{theorem}
\label{thm:unbiased}
Condition on $\mathcal F_t$. Assume (i) all policies, prompt laws, nuisance
models, counts, and $\alpha$ are fixed before sampling; (ii)
\cref{eq:target-law} is differentiable and differentiation and expectation may
be interchanged; (iii) the deterministic reward does not directly depend on
$\theta$; (iv) $b_t$ depends only on $X$; (v) current and stale samples are
conditionally independent, newly drawn, and have not influenced $\theta_t$;
(vi) $P_t$ and $G_t$ have exactly evaluable masses on a common sample space, so
$M_t$ is computed exactly; (vii) every preassigned draw reaches a scored
outcome almost surely; and (viii) required moments are finite. With exact
\cref{eq:weight} and none of the clipping or normalization operations excluded
above,
\begin{equation}
\mathbb E[\widehat g_t\mid\mathcal F_t]
=\nabla_\theta\mathcal L(\theta_t,\lambda_t).
\end{equation}
\end{theorem}

The proof expands the two strata:
\begin{align}
\mathbb E[\widehat g_t]
&=\mathbb E_P[(h-b)s]
+\alpha\mathbb E_P[w(U-h)s]\nonumber\\
&\quad +(1-\alpha)\mathbb E_G[w(U-h)s]\\
&=\mathbb E_P[(h-b)s]+\mathbb E_M[w(U-h)s]\nonumber\\
&=\mathbb E_P[(U-b)s]=\mathbb E_P[Us].
\end{align}
The last equality is the prompt-baseline identity.

If $\widetilde w(z)$ is the same fixed, pointwise implemented weight in both
strata and
$m_0(z)=\mathbb E[U\mid Z=z]$, the population bias is
\begin{equation}
\mathbb E_M[(\widetilde w-w)(m_0-h)s].
\label{eq:bias}
\end{equation}
This is a weight-misspecification identity only under frozen common weights.
For source-specific fixed implementations it becomes
\begin{equation}
\alpha\mathbb E_P[(\widetilde w_P-w)(m_0-h)s]+
(1-\alpha)\mathbb E_G[(\widetilde w_G-w)(m_0-h)s].
\end{equation}
Thus either exact density ratios or an exact action-conditional outcome model
removes this bias. Low prediction error or calibration alone is insufficient,
and focal-batch nuisance fitting, batch-normalized weights, and data-dependent
weights are not covered. In
deterministic RLVR, $m_0(z)=U(z)$, so an exact outcome model makes the stale
residual vanish and leaves an on-policy direct estimator. We therefore
describe \cref{eq:estimator} as augmented importance sampling rather than
relying on a broad ``doubly robust replay'' label.

\paragraph{Scope of the result.}
\Cref{thm:unbiased} applies before one update. After the policy changes, the
current component of $M_t$ is no longer the target policy, the
$1/\alpha$ bound need not hold, and the same data influenced the updated
policy. Multiple epochs, ratio recomputation, or small KL do not restore the
theorem. Optimizer momentum and clipping may still be useful, but the theorem
is about the raw estimator rather than the transformed displacement or
monotonic improvement.

\section{Experiments}
\label{sec:experiments}

Evaluation starts with a deterministic software Gate 0 on a finite trajectory
law. It must exactly reproduce the target gradient, mixture normalization,
weight bound, and analytic bias of each injected ratio error before Phase A.
Phase A then establishes estimator behavior at a frozen language-model policy
before Phase B makes training claims. The primary hypotheses are:
\textbf{H1}, exact adaptive-mixture augmentation reduces projected gradient MSE
relative to on-policy estimation at matched total compute; \textbf{H2}, any
reduction translates to better held-out reward per generated token under
one-update-per-batch training; and \textbf{H3}, the expected-response-length constraint
produces a better reward--length frontier than implicit response-length
normalization. Each can fail independently.

\subsection{Experimental Setup}

\paragraph{Frozen-policy audit.}
For each of two model families and both math and code verifiers, freeze a
current policy, one stale policy, nuisance models, prompts, and decoders.
For each target law generate two mutually independent large $P$-only reference
estimates $R_1,R_2$, independent also of candidate estimate $\widehat g$, on
adapter layers and projections frozen before outcomes. We estimate projected
MSE by
\begin{equation}
\widehat{\operatorname{MSE}}_\times=
(\widehat g-R_1)^\top(\widehat g-R_2).
\end{equation}
Its expectation is projected MSE although one realization may be negative.
Intervals independently resample candidate replications, $R_1$ trajectories,
and $R_2$ trajectories, retaining their shared-reference dependence.
Selection and confirmation use separate reference pairs under their respective
laws. Nuisance and proposal
models are trained on disjoint data. Repeated subsamples compare on-policy
policy gradient, ordinary
importance sampling, bounded mixture importance sampling, augmentation with
several predictor qualities, and adaptive versus uniform $Q_t$. Endpoints are
bias, variance, MSE, interval coverage, effective sample size (ESS), and
compute. Fault injection deliberately omits prompt ratios, uses base-model
rather than decoder probabilities, fits the focal nuisance outcome, clips
weights, and reuses data that formed the target policy.
The committed finite-law Gate~0 currently executes the omitted-prompt-ratio
and clipping cases. Base-versus-decoder scoring, focal nuisance fitting, and
reuse are pending executable validators and must pass before Phase~A; listing
them here is a test plan, not evidence that they have run.
The committed public-data configuration is a one-family mathematics pre-pilot
for H1 only. It uses disjoint normalized laws $D_{\rm fit}$,
$D_{\rm select}$, and $D_{\rm confirm}$; selection MSE is not described as an
estimate under the confirmation or full-frame law. For
$\alpha\in\{.25,.5,.75,1\}$ it fixes
$(n_P,n_G)\in\{(16,48),(32,32),(48,16),(64,0)\}$ and minimizes projected MSE
times measured accelerator-seconds at a common 64-trajectory and verifier-call
budget. Tying mixture mass to allocation fraction is a pilot restriction, not
a condition for unbiasedness; a confirmatory design must vary them separately.
Code, second-family, H2, and H3
confirmatory runs require a later frozen configuration.

\paragraph{End-to-end training.}
Confirmatory training uses two open model families at small and medium scale,
one exact-answer math suite and one sandboxed code suite. Prompts, hidden
verifiers, maximum generation length, temperature, optimizer, and checkpoint
selection rules are frozen before the main runs. Each collected batch produces
one update and is never used again by the theorem-facing method. Staleness is
stratified by measured policy KL, not checkpoint age.

Baselines are on-policy RLOO, Dr.\ GRPO, separately tuned DAPO and GSPO, TOPR,
variance-informed prompt allocation (VIP), DBB reward estimation,
PrefixRL, PROS, SPEC-RL, MinPRO,
the ICML replay design, POPO,
plain bounded-mixture importance sampling, augmented sampling with uniform
$Q$, and the full adaptive proposal. All generated prefixes and complete
sequences, verifier calls, likelihood passes, nuisance training, policy
updates, and hyperparameter searches are charged. Each configuration receives
the same token and verifier-call budget, while total allocated accelerator and
CPU time is reported as a Pareto cost rather than assumed equal. We use five
independent development seeds and five frozen confirmatory seeds.
Source-isolation controls set $\mu_t=\pi_t$ while adapting $Q_t$, and set
$Q_t=D$ while varying policy staleness. Rollout counts are fixed before each
batch. A batch is admitted only if its predeclared worst-case token and
verifier-call reserve fits; the run stops at the last fully admitted batch and
reports budget undershoot. Realized token, call, accelerator, and CPU costs are
never converted into one unspecified scalar. Reward-versus-realized-token and
reward-versus-accelerator curves are primary; endpoint contrasts are reported
only at their largest common realized-cost checkpoint, with no interpolation
across updates.

\subsection{Main Results}

\begin{table*}[t]
\caption{Primary matched-budget results. Higher held-out reward is better;
lower projected gradient MSE and generated tokens to threshold are better.
Dashes mark measurements that have not yet been run.}
\label{tab:main}
\centering
\begin{small}
\begin{tabular}{lcccc}
\toprule
Method & Gradient MSE & Held-out reward & Tokens to threshold
& Mean response tokens \\
\midrule
On-policy RLOO & -- & -- & -- & -- \\
Bounded mixture IS & -- & -- & -- & -- \\
Augmented mixture, uniform $Q$ & -- & -- & -- & -- \\
Augmented mixture, adaptive $Q$ & -- & -- & -- & -- \\
Practical repeated-replay variant & -- & -- & -- & -- \\
\bottomrule
\end{tabular}
\end{small}
\end{table*}

Every value in \cref{tab:main} will aggregate independent runs with a
predeclared interval and link to immutable trajectory manifests. The
repeated-replay row is explicitly not covered by \cref{thm:unbiased}; it tests
whether a biased engineering variant is nevertheless more useful.

\subsection{Ablations and Diagnostics}

We ablate the on-policy anchor fraction $\alpha$, augmentation, adaptive prompt
proposal, prompt baseline, cross-iteration nuisance freezing, and length dual.
A separate deliberate-violation study runs 1, 2, 4, and 8 epochs; uses fresh,
previously consumed, and reward-selected stale data; and compares unclipped and
clipped ratios. After each update, a large fresh on-policy batch estimates the
resulting approximation error.

Required diagnostics include the numerical $1/\alpha$ bound; mixture
normalization; joint and source-specific ESS; maximum weight and support
failures; direct and residual gradient norms and cosine; outcome calibration by
length, difficulty, and source; fresh versus augmented projected gradients;
pre/post-update KL; entropy; sign-conditioned update mass; policy-to-behavior
KL; dual trajectory; reward--token frontier; parser errors; and verifier
disagreement. Reporting only reward curves is insufficient to test the method.
The length study freezes
$\kappa\in\{512,1024,2048\}$, $\lambda_0=0$, and dual step size $10^{-4}$;
last-iterate and final-20\%-average frontiers are both reported. Replay
comparators freeze admission, buffer capacity, reuse count, KL stratum,
clipping, and update count, with original generation and rescoring charged.

\subsection{Generalization and Safety}

Generalization holds the estimator configuration fixed across model family,
scale, prompt difficulty, context length, and verifier domain. Hidden verifier
variants alter answer formatting, unit tests, parser implementation, and
timeouts. Reward controls include shuffled labels, random rewards, weak
heuristics, and anti-correlated signals. We measure train/hidden reward gaps,
pass@1 and pass@$k$, response-length quantiles, invalid-output rates,
memorized-answer indicators, and diversity. Safety claims are not inferred
from verifier reward; the protocol will report separate capability and
harmful-behavior audits as possible regressions.

\section{Limitations}
\label{sec:limitations}

The assumptions intentionally exclude the main appeal of replay: repeated use
of old data. Single-use stale actors may add systems complexity without enough
variance reduction to justify it. Exact sequence probabilities are expensive
and fragile when production decoders use top-$p$, rejection, tool calls, or
environment transitions. Adaptive prompt density must be logged exactly.
Importance weights remain high-variance despite their bound, especially for
small $\alpha$ and long horizons.

The outcome predictor is a critic by another name and may consume substantial
compute. The two-route property in \cref{eq:bias} is practically asymmetric:
behavior likelihoods are logged, whereas an exact sequence-conditional verifier
predictor is unrealistic. Conditional gradient unbiasedness does not imply
low MSE, stable Adam updates, monotonic improvement, or resistance to reward
hacking. Finally, math and code evidence may not generalize to subjective
human-feedback rewards or frontier models.

\section{Conclusion}
\label{sec:conclusion}

We separated statistically controlled stale generation from ordinary replay
and specified an adaptive-mixture augmented policy gradient for
expected-response-length-constrained RLVR. The estimator has a simple
conditional expectation identity under
strong, auditable restrictions, and fails to inherit that identity after data
reuse or multiple updates. This clarity is useful only if the estimator
actually reduces gradient error after all added compute is charged. The
frozen-policy audit is therefore the next decisive step; downstream training
should proceed only if that test succeeds.

% Acknowledgements must not appear in the anonymous review submission.
% \section*{Acknowledgements}

\section*{Impact Statement}

More sample-efficient RL could reduce the compute required to improve
reasoning systems, but it can equally accelerate optimization toward flawed or
harmful verifiers. Exact-looking tests can encode incomplete specifications,
reward insecure code, or favor brittle answer formats. Adaptive allocation may
also neglect rare prompts unless its support floor is enforced. We therefore
require hidden-verifier variants, random and adversarial reward controls,
minimum prompt support, full token and energy accounting, and separate safety
audits. Stale trajectories and prompts may contain sensitive content and must
follow the same retention and access controls as training data. The paper does
not equate verifiability with correctness, harmlessness, or social value.

\bibliography{references}
\bibliographystyle{../template/icml2026/icml2026}

\clearpage
\appendix
\onecolumn

\section{Algorithmic and Reproducibility Details}
\label{app:algorithm}

\subsection{Dual Update}
The expected-response-length dual uses only fresh target-policy trajectories:
\begin{equation}
\lambda_{t+1}=
\left[\lambda_t+\eta_\lambda
\left\{\frac1{n_P}\sum_{i=1}^{n_P}C(Y_i^P)-\kappa\right\}\right]_+.
\end{equation}
$\kappa$ is fixed before training. Policy and dual updates both estimate
quantities at $(\theta_t,\lambda_t)$ and are committed after collection.
Reward, response length, and violation are reported separately rather than
only through $U_\lambda$. This projected stochastic controller does not
guarantee that every iterate or the final policy is feasible. The constraint
does not cover prompt tokens,
likelihood passes, verifier cost, or total training compute; those belong to
the experimental ledger.

\subsection{Logging Contract}
Each trajectory records prompt identifier and probability, source component,
current and behavior checkpoint hashes, tokenizer, temperature, EOS and
truncation rule, random seed, exact current and behavior sequence
log-probabilities, reward components, parser status, generated tokens, verifier
version, nuisance-model revision, use count, and policy-update identifier. The
run aborts if a supposedly theorem-facing trajectory has use count above one
or if its source probability cannot be reconstructed.
For all candidate trajectories, including those in aborted batches, the
ledger computes
\begin{align}
K_{\mathrm{gen}}&=\sum_z C(z),&
K_{\mathrm{ver}}&=\sum_z\mathbf 1\{\text{verifier called on }z\},\\
K_{\mathrm{score}}&=\sum_z C(z)\,N_{\mathrm{added\ score}}(z).
\end{align}
No cost disappears when a trajectory is excluded from the estimator.

\subsection{Protocol Manifest}
The repository runner
\texttt{python3 scripts/run\_experiment\_pipeline.py
--plan experiments/rl\_pipeline.json} enumerates immutable prompt acquisition,
verifier processing, policy and nuisance-model construction, and a frozen-policy
H1 estimator pre-pilot. The full end-to-end H2/H3 plan is not executable from
this configuration. The committed JSON is a
stage and artifact contract tied to the hashed RL section of
\texttt{experiments/pilot\_preregistration.json}; site-local commands are
required to execute it, and each executed stage records output hashes and
separate token, verifier, accelerator, CPU, and wall-time costs. The public
source registry pins the prompt source and base checkpoint only. The current
and stale policy tensor hashes, tokenizer/decoder hashes, KL sampling law, and
mathematics-verifier revision are still unfrozen blockers; this H1 pre-pilot
cannot execute until a replacement configuration records them. A hidden code
evaluator is outside this mathematics-only pre-pilot and must be registered by
any later code study.
Automated license metadata is only a preflight: real acquisition remains
blocked on manual provenance, terms, privacy, and leakage review.
Once approved, \texttt{scripts/materialize\_public\_data.py} verifies each
site-local export hash, emits row-level provenance, deduplicates within roles,
and quarantines normalized-content collisions across roles. The committed
fixture exercises those rules without downloading prompts; all of its outputs
are marked non-empirical.
\texttt{scripts/run\_content\_gates.py} also binds the frozen
\texttt{experiments/content\_gates.json} contract to the materialization hashes
and fails closed on retained named PII/secret patterns or lexical
near-duplicates. Because the large-artifact MinHash-LSH candidate screen can
miss pairs, it is always inconclusive and blocking until a reviewed
recall-complete audit is hash-bound. Even an exhaustive pass neither clears a
source nor establishes semantic independence from a hidden evaluator.
The companion \texttt{synthetic\_commands.json} fixes a four-cell trajectory
law, $\alpha=1/2$, 32 draws per source, 2,000 replications, and seed 20260902.
It enumerates the target and estimator expectations to machine precision,
checks $\sum_zM(z)w(z)=1$ and $w\le1/\alpha$, and verifies the bias identity
for omitted-prompt-ratio and weight-clipping errors. Gate 0 fails closed if any
identity differs by more than $10^{-12}$. This is an equation and software
audit, not model-training evidence; its outputs cannot populate
\cref{tab:main}.

\subsection{Optional Prefix Gate}
Prefix continuation is excluded from the core experiments. A later extension
may draw a prefix $H\sim\mu_t$, choose a predictable
$p_t(H)\in[p_{\min},1]$ before observing the suffix or reward, and include a
completed stale-source residual with factor $I/p_t(H)$. The denominator remains
the fixed number of candidate prefixes, not the number completed. Every
discarded prefix is charged. Sequential gates require the product of all
conditional continuation probabilities; group-relative estimators additionally
require joint inclusion probabilities. This extension must be compared
directly with PrefixRL and PROS \cite{setlur2026prefixrl,huang2026pros}.

\section{Additional Results and Failures}
\label{app:additional}

\begin{table}[h]
\caption{Mandatory estimator-audit diagnostics; pending execution.}
\centering
\begin{tabular}{lc}
\toprule
Diagnostic & Value \\
\midrule
Projected bias relative to fresh target & -- \\
Projected variance and MSE & -- \\
Maximum mixture weight / $1/\alpha$ & -- \\
Joint effective sample size & -- \\
Direct--residual gradient cosine & -- \\
Prompt-density reconstruction error & -- \\
Outcome calibration error & -- \\
Added compute / on-policy compute & -- \\
\bottomrule
\end{tabular}
\end{table}

Complete learning curves, per-seed endpoints, theorem-violation results,
reward-control trajectories, parser failures, and examples selected by a
fixed random seed will be released whether they favor the method or not.

\end{document}
